\documentclass{article}
\usepackage[margin=1in]{geometry}
\usepackage{hyperref}

\title{Problem Set 4}
\author{Yeyang Zhou}
\date{\today}

\begin{document}
\maketitle

\section{Data Sources of Interest}

I am interested in scraping data from several sources. First, SEC EDGAR provides 
freely accessible 10-K and 10-Q filings, which are valuable for empirical accounting 
research on corporate disclosures. Second, the Federal Reserve Economic Data (FRED) 
API offers macroeconomic time series data including GDP, inflation, and interest rates, 
which are useful for finance research. Third, RavenPack and similar news analytics 
platforms provide event-driven datasets that can be used to study market reactions 
to corporate announcements such as layoffs or earnings surprises.

\section{Exercise 5 Answers}

\subsection{Class of mydf and mydf\$date}
The object \texttt{mydf} is of class \texttt{tbl\_df}, \texttt{tbl}, and 
\texttt{data.frame}. The column \texttt{mydf\$date} is of class \texttt{character}. 
This is because JSON stores all values as strings by default, so the date column 
is not automatically parsed into a native \texttt{Date} object.

\subsection{First 6 Rows of mydf}
The first 6 rows of \texttt{mydf} all have date value of 1 (year 1 AD), with 
descriptions of historical events from the Roman Empire and Asia, along with 
language, category, and granularity information.

\section{Exercise 6 Answers}

\subsection{Class of df1 and df}
\texttt{df1} is of class \texttt{tbl\_df}, \texttt{tbl}, and \texttt{data.frame}, 
stored in R memory. \texttt{df} would be of class \texttt{tbl\_spark}, which is 
a reference to a distributed dataset in Spark rather than a local R object.

\subsection{Column Name Differences}
The column names differ between the two objects. R allows dots in column names 
(e.g., \texttt{Sepal.Length}), but Spark and SQL do not support dots in column names. 
Therefore, when copying to Spark, \texttt{sparklyr} automatically converts dots 
to underscores (e.g., \texttt{Sepal\_Length}).

\subsection{Note on sparklyr Installation}
The \texttt{sparklyr} package was not available for R version 4.0.2 on OSCER, 
so the Spark exercises could not be executed. The code is written in 
\texttt{PS4b\_Zhou.R} based on the assignment instructions.

\end{document}